% ================================================================
%  chapters/parts/ch06-a.tex
%  第 6 章  催化过程的拉曼光谱方法（分片 a）
%  源：S1（上册）第 6 章，书内 p.343—372 / PDF p.362—391
%  覆盖：6.1 拉曼光谱原理简述；6.2 拉曼光谱实验技术的发展
%  分片 b 覆盖 6.3 及以后（催化应用、紫外拉曼、SERS、共焦）
% ================================================================

\chapterstart{催化过程的拉曼光谱方法}%
{本章导读：拉曼光谱测的是\textbf{分子极化率变化}，与红外光谱测偶极矩变化恰成互补，
因此它特别擅长看骨架振动、金属氧化物和负载物种，而且几乎不受水的干扰。本章回答三段问题：
拉曼散射的信号从哪个物理量来、为什么同一分子的拉曼位移不随激发波长改变；荧光为什么会
把拉曼信号整个淹没、有哪几条可行的抑制路线；以及一台拉曼光谱仪的五个部件各自把什么
条件写进了谱图。6.1 节讲原理，6.2 节压缩为部件对照表。催化应用、紫外拉曼、表面增强与
共焦成像见后续各节。}
\label{ch:06}

\subsection{拉曼光谱原理简述}

\subsubsection{拉曼效应：测什么、为什么值得测}

\begin{probout}
拉曼光谱的\textbf{输入}是一束频率为 $\nu_0$ 的单色激光和一块样品，\textbf{输出}是散
射光强度对\textbf{拉曼位移}的曲线。它真正测量的物理量只有一个：\textbf{分子极化率张量
对简正坐标的导数} $\left(\partial \alpha / \partial q\right)_0$。这个导数是否为零，决
定了一条振动是「拉曼活性」还是「拉曼非活性」；它的大小决定峰有多强。因此拉曼给出的
是\textbf{分子骨架与化学键的振动信息}，配合谱峰位置可作物相识别，配合谱峰位移可作物
种与载体相互作用的判据。它不直接给出元素组成，也不直接给出表面覆盖度；强度只能相对
比较（见§6.2.1 的「数据到结论」）。
\end{probout}

印度物理学家 C.\,V.\,Raman 于 1928 年用汞灯某一波长照射苯、\ce{CCl4} 等液体，将散射
光经分光棱镜后用眼睛或摄谱仪观察，发现散射光中除与入射光频率 $\nu_0$ 相同的强瑞利
（Rayleigh）散射外，两侧还对称地分布着一系列其他频率的非弹性散射光，其强度比瑞利散
射弱得多，通常只有瑞利光强度的 \num{e-9}—\num{e-6} 量级。这一现象称为\textbf{拉曼
散射效应}，他因此获得 1930 年诺贝尔物理学奖\src{S1 §6.1.1，书内 p.345 / PDF p.364}。

拉曼谱线的频率随入射光频率改变，但\textbf{拉曼散射光与瑞利散射光的频率差基本不随入射
光频率变化}，只与样品分子的振动和转动能级有关。频率大于瑞利线的称反斯托克斯光
（anti-Stokes），小于瑞利线的称斯托克斯光（Stokes）；瑞利线与斯托克斯线的频率差即
\textbf{拉曼位移}。常温下斯托克斯线比反斯托克斯线强几个数量级。拉曼谱线强度与入射光
强度、样品分子浓度成正比，据此可作结构与浓度分析。

\begin{principle}
\textbf{拉曼位移的定义}。设入射光频率为 $\nu_0$、拉曼散射光频率为 $\nu_{\mathrm{s}}$，
分子振动频率为 $\nu_{\mathrm{m}}$，则
\begin{equation}
\begin{aligned}
\Delta \tilde\nu \;&=\; \frac{\nu_0-\nu_{\mathrm{s}}}{c} \;=\; \frac{\nu_{\mathrm{m}}}{c}\\[1pt]
&\text{（斯托克斯线：}\ \nu_{\mathrm{s}}=\nu_0-\nu_{\mathrm{m}}\text{）}
\end{aligned}
\label{eq:raman-shift}
\end{equation}
\textbf{成立条件}：单光子非弹性散射、分子处于电子基态、无多光子过程。用波数
$\tilde\nu=\nu/c$ 表述时，$\Delta\tilde\nu$ 的单位是 \unit{cm^{-1}}，与激发波长无关——
这正是拉曼位移被称为「分子本征物理量」的原因。\textbf{量纲检查}：分子与分母同除光速，
$\Delta\tilde\nu$ 与 $\nu_{\mathrm{m}}/c$ 同为波数量纲（\unit{m^{-1}}，习惯写
\unit{cm^{-1}}），无误。
\end{principle}

拉曼之所以在催化研究中有位置，来自以下几条互相独立的性质\src{S1 §6.1.1、§6.3，
书内 p.345、p.355 / PDF p.364、p.374}：

\begin{itemize}
\item \textbf{给的是催化剂本身及表面上物种的结构信息}，这是认识催化剂与催化反应最直接
  的一类信息；
\item \textbf{容易实现原位条件下的研究}（高温、高压、复杂体系），气相光谱干扰很弱；
\item \textbf{水的拉曼散射很弱}，因而比红外更适合水相体系，特别适于跟踪催化剂制备从
  水相到固相的实时过程；
\item \textbf{可测低波数}。红外一般很难得到 \qty{200}{cm^{-1}} 以下的光谱，拉曼可以测
  到几十个波数，而低波数区正反映骨架与晶格结构；
\item \textbf{常用载体（\ce{gamma-Al2O3}、\ce{SiO2} 等）的拉曼散射截面很小}，对表面
  负载物种的干扰少。
\end{itemize}

\begin{pitfall}
上述优点都有边界，不能当成无条件成立。催化研究中\textbf{荧光干扰和灵敏度相对较低仍是
拉曼原位研究的难题}：载体本身（尤其是氧化铝）常带强荧光，使低负载量氧化物样品几乎测不
出来。另外，低波数优势只在能有效抑制瑞利线的多联光谱仪上才兑现，单联光谱仪的杂散光足
以盖掉 \qty{200}{cm^{-1}} 以下的弱峰。
\end{pitfall}

\subsubsection{拉曼与红外的互补性：互斥规则}

一条振动是否可见，取决于它调制的是哪个物理量\textbf{：红外要求分子的偶极矩发生变化
（$\partial \mu/\partial q \neq 0$），拉曼要求分子的极化率发生变化
（$\partial \alpha/\partial q \neq 0$）。}二者并不等价，因此两种光谱\textbf{有一定程度的
互补性，而不能互相代替}\src{S1 §6.3，书内 p.355 / PDF p.374}。对\textbf{有对称中心的
分子}，互斥规则（rule of mutual exclusion）进一步收紧：凡红外活性的振动拉曼非活性，
反之亦然。这一条对催化最有用的推论是——\textbf{全对称伸缩振动（骨架呼吸、端基
\ce{M{=}O} 伸缩）几乎只在拉曼里强，而极性基团的强吸收（\ce{OH}、\ce{NH}、\ce{CO}）几乎
只在红外里强}。表~\ref{tab:raman-ir} 把两者的分工与局限一次列清。

\begin{table}[htbp]
\centering
\caption{拉曼与红外光谱的互补性对照。读者应看「选律」与「样品适应性」两行：前者决定
两者给出的是不同子集的振动，后者决定哪一种在含水、含载体、低波数场景下更可用。
局限一行提示两种方法各自最容易失败的地方。
\src{据 S1 §6.3，书内 p.355 / PDF p.374 归纳。}}
\label{tab:raman-ir}
{\srcsize
\begin{tabularx}{\columnwidth}{@{}lXX@{}}
\toprule
 & 拉曼光谱 & 红外光谱 \\
\midrule
测什么 & 极化率变化 $\partial\alpha/\partial q$ & 偶极矩变化 $\partial\mu/\partial q$ \\
选律 & 有对称中心时与红外互斥；全对称振动强 & 有对称中心时与拉曼互斥；极性基团强 \\
最擅长的振动 & 骨架振动、晶格振动、\ce{M-O}、\ce{M{=}O}、\ce{S-S}、\ce{C{=}C} & \ce{OH}、\ce{NH}、\ce{CO}、\ce{NO} 等极性基团、吸附探针分子 \\
样品适应性 & 水干扰弱，可测水相与原位气相；可直接测粉末，无需压片 & 水干扰强，需干燥与压片/稀释；水相体系困难 \\
低波数能力 & 可低至几十 \unit{cm^{-1}}（需多联光谱仪抑瑞利线） & 通常难低于 \qty{200}{cm^{-1}} \\
载体影响 & \ce{gamma-Al2O3}、\ce{SiO2} 等散射截面小，干扰少 & \qty{1000}{cm^{-1}} 以下载体吸收强，几乎不透 \\
主要局限 & 荧光干扰、散射截面小、信号弱、无绝对强度 & 载体强吸收、需样品制备、高温黑体辐射 \\
\bottomrule
\end{tabularx}}
\end{table}

\emph{核对说明}：已对照原页面图像（书内 p.355 / PDF p.374）。原书 §6.3
开头即给出本表「选律」行的依据：「拉曼光谱与红外光谱都能得到分子振动和转动光谱，
但分子的极化率发生变化时才能产生拉曼活性，对于红外光谱，只有分子的偶极矩发生变化时才具有红外活性，
因此二者有一定程度的互补性，而不可以互相代替」；「低波数能力」行同样出自该页：「红外光谱一般很难得到低波数（200
cm$^{-1}$ 以下）的光谱，但拉曼光谱甚至可以得到几十个波数的光谱」；「载体影响」行亦在该页：
「由于常用载体（如 $\gamma$-Al$_2$O$_3$ 和 SiO$_2$ 等）的拉曼散射截面很小，
因此载体对表面负载物种的拉曼光谱的干扰很少。而大部分载体（如 $\gamma$-Al$_2$O$_3$、
SiO$_2$ 和 TiO$_2$ 等）在低波数的红外吸收很强，在 1000 cm$^{-1}$
以下几乎不透过红外光」；「样品适应性」中的水相一句同样见于该页：「由于水的拉曼散射很弱，
因此拉曼比红外更适合进行水相体系的研究」。表中互斥规则、「最擅长的振动」一列、「可直接测粉末、无需压片」
以及「主要局限」一列（其中「荧光干扰和灵敏度较差是阻碍常规拉曼光谱得到广泛应用的最大问题」一句另见书内 p.374
/ PDF p.393）属通用方法学归纳，原书该处未逐条给出。


\begin{crosscheck}
两句话的实用判据：\textbf{要问「骨架变没变」用拉曼，要问「羟基和酸性位怎么样」用红外}。
对同一个样品，\ce{MoO3/Al2O3} 的 \ce{Mo{=}O} 端基伸缩振动与 \ce{Mo-O-Mo} 桥键在拉曼里
清晰可辨，而载体 \ce{gamma-Al2O3} 在 \qty{1000}{cm^{-1}} 以下的红外几乎不透光，红外在
这一段基本没有可用窗口。反过来，吡啶吸附探针识别 Br\o nsted 酸与 Lewis 酸是红外的经
典做法，拉曼可作补充——即使在 \qty{1200}{cm^{-1}} 以下，拉曼仍能得到表面物种的光谱。
回查的规则是：\textbf{任何「某物种存在」的结论，若只有一条谱峰支撑，都应该用另一条谱
峰、另一种探针分子或另一种表征手段交叉确认}。
\end{crosscheck}

\subsubsection{拉曼散射的经典理论}

按电磁辐射的经典图像，频率为 $\nu_0$ 的入射电磁场使分子电子云发生变化，感生出振荡的
电磁多极子，振荡多极子又向外辐射。辐射频率与入射频率相同为瑞利散射；辐射频率为
$\nu_0\pm\nu_{\mathrm{m}}$ 则为拉曼散射。因此\textbf{拉曼效应可归结为：外电场感生的
电偶极矩被分子的振动运动所调制}。设入射光为单色光

\begin{equation}
E=E_0\cos 2\pi\nu_0 t ,
\label{eq:incident-field}
\end{equation}

入射光不太强时，感生电偶极矩 $\mathbf{p}$ 可近似写成分子极化率张量
$\boldsymbol{\alpha}$ 与场强的乘积

\begin{equation}
\mathbf{p}=\boldsymbol{\alpha}\mathbf{E}, \qquad
\mathbf{p}=\boldsymbol{\alpha}\mathbf{E}_0\cos 2\pi\nu_0 t .
\label{eq:induced-dipole}
\end{equation}

分子中各原子核在平衡位置附近振动，极化率随之改变。把极化率张量的每个分量在平衡位置
处按简正坐标 $q$（振动频率 $\nu_{\mathrm{m}}$）作 Taylor 展开：

\begin{equation}
\alpha=\alpha_0+\left(\frac{\partial\alpha}{\partial q}\right)_{\!0} q
+\frac12\left(\frac{\partial^2\alpha}{\partial q^2}\right)_{\!0} q^2+\cdots,
\qquad q=q_0\cos 2\pi\nu_{\mathrm{m}} t .
\label{eq:alpha-expansion}
\end{equation}

把式~\eqref{eq:alpha-expansion} 代回式~\eqref{eq:induced-dipole}，用
$\cos A\cos B=\tfrac12[\cos(A+B)+\cos(A-B)]$ 展开，得到三项——即瑞利散射项、反斯托克斯
项与斯托克斯项：

\begin{align}
p_{\text{瑞利}} &= \alpha_0 E_0\cos 2\pi\nu_0 t, \nonumber\\[1pt]
p_{\text{反斯托克斯}} &= \tfrac12
  \left(\frac{\partial\alpha}{\partial q}\right)_{\!0} q_0 E_0
  \cos 2\pi(\nu_0+\nu_{\mathrm{m}})t, \nonumber\\[1pt]
p_{\text{斯托克斯}} &= \tfrac12
  \left(\frac{\partial\alpha}{\partial q}\right)_{\!0} q_0 E_0
  \cos 2\pi(\nu_0-\nu_{\mathrm{m}})t .
\label{eq:three-terms}
\end{align}

\textbf{由此得到拉曼活性的经典判据}：$\left(\partial\alpha/\partial q\right)_0\neq 0$ 才
有拉曼活性，等于零则为非拉曼活性——这就是「拉曼测极化率变化」的确切含义
\src{S1 §6.1.2，书内 p.347—348 / PDF p.366—367}。振动的感生电偶极所辐射的散射功率为

\begin{equation}
P_{\mathrm{sc}}=E_0\,\frac{\mathrm{d}\sigma}{\mathrm{d}\Omega}\,\Delta\Omega ,
\qquad E_0=\frac{P_0}{A_{\mathrm{beam}}} ,
\label{eq:classical-intensity}
\end{equation}

式中 $\theta$ 为观察方向与感生偶极矩的夹角。\textbf{先把五个量的口径说清}（全部用 SI
单位）：$P_0$ 是入射到样品上的\emph{功率}（\unit{W}）；$A_{\mathrm{beam}}$ 是光束在样品
处的\emph{光斑面积}（\unit{m^2}）；$E_0=P_0/A_{\mathrm{beam}}$ 是样品处的\emph{入射辐照度}
（\unit{W.m^{-2}}）；$\mathrm{d}\sigma/\mathrm{d}\Omega$ 是\emph{微分散射截面}
（\unit{m^2.sr^{-1}}），按定义即「单位立体角内的散射功率除以入射辐照度」，\textbf{它本身
就含所观察方向的角分布}——对感生电偶极辐射，
\begin{equation}
\frac{\mathrm{d}\sigma}{\mathrm{d}\Omega}=K\sin^2\theta ,
\label{eq:dipole-dcs}
\end{equation}
其中 $K\equiv\left.(\mathrm{d}\sigma/\mathrm{d}\Omega)\right|_{\theta=90^\circ}$ 是与方向
无关的\emph{参考方向系数}（偶极垂直方向的微分截面），它对应原书式中的前因子
$k'_{\omega}\omega_{\mathrm{s}}^{4}p^{2}$；
$\Delta\Omega$ 是接收端张开的\emph{接收立体角}（\unit{sr}），由接收面积
$A_{\mathrm{rec}}$ 与样品到接收端的距离 $R$ 决定，$\Delta\Omega\approx A_{\mathrm{rec}}/R^2$。
把式~\eqref{eq:dipole-dcs} 代回式~\eqref{eq:classical-intensity}，\textbf{偶极的角分布
$\sin^2\theta$ 只经微分散面出现一次}，落入接收端的散射功率为
$P_{\mathrm{sc}}=E_0K\sin^2\theta\,\Delta\Omega$。若只问\emph{比例}，两端同除 $P_0$：

\begin{equation}
\frac{P_{\mathrm{sc}}}{P_0}
=\frac{\mathrm{d}\sigma}{\mathrm{d}\Omega}\,\frac{\Delta\Omega}{A_{\mathrm{beam}}}
=\frac{K\sin^2\theta\,\Delta\Omega}{A_{\mathrm{beam}}} ,
\label{eq:classical-intensity-ratio}
\end{equation}

\textbf{量纲检查}：左端是功率之比，无量纲；右端为
$\unit{m^2.sr^{-1}}\times\unit{sr}\times\unit{m^{-2}}$（$K\sin^2\theta$ 无量纲）——立体角
与面积都相消，也是无量纲，两端一致。分母上的 $A_{\mathrm{beam}}$ 不能省：它正是「入射功率 → 入射
辐照度」那一步带进来的因子；上一轮把该式写成
$I/I_0=(\mathrm{d}\sigma/\mathrm{d}\Omega)\,\Delta\Omega\,\sin^2\theta$，右端只剩面积量纲
（\unit{m^2}），与无量纲的左端不一致，本稿据此改写为
式~\eqref{eq:classical-intensity-ratio}。

\textbf{观测几何与四条可检验的推论}：偶极辐射的强度按 $\sin^2\theta$ 分布——在与偶极矩
垂直的方向（$\theta=90^\circ$）最强，沿偶极矩方向（$\theta=0^\circ$）为零。以 90° 方向为
基准，45° 与 30° 的相对收集功率分别是 $\sin^2 45^\circ/\sin^2 90^\circ=0.5$ 与 $0.25$。
在角分布不变的前提下，式~\eqref{eq:classical-intensity} 给出四条可以逐条代入核对的结论：

\begin{enumerate}
\item 固定入射功率与光斑，只把接收立体角加倍（接收面积加倍，或距离减为 $1/\sqrt2$）：
      $P_{\mathrm{sc}}$ 加倍；
\item 固定 $\Delta\Omega$（距离拉长时同步放大接收面积以保持立体角不变），距离 $R$ 加倍：
      $P_{\mathrm{sc}}$ \emph{不变}——微分散射截面本身就定义在「单位立体角」上，几何衰减
      已由 $\Delta\Omega$ 承担，再乘一次 $R^{-2}$ 就是重复计入；
\item 固定接收\emph{面积} $A_{\mathrm{rec}}$，距离 $R$ 加倍：$\Delta\Omega$ 降为 $1/4$，
      故 $P_{\mathrm{sc}}$ 约为原来的 $1/4$；
\item 比值 $P_{\mathrm{sc}}/(P_0\Delta\Omega)
      =K\sin^2\theta/A_{\mathrm{beam}}
      =(\mathrm{d}\sigma/\mathrm{d}\Omega)/A_{\mathrm{beam}}$ 与 $\Delta\Omega$
      无关——这正是「比较散射截面时必须同时固定收集立体角与光斑面积」的原因。
\end{enumerate}

以上四条只针对\emph{单个散射体}、固定入射辐照度与固定角分布范围的情形：$N$ 个相互独立
的散射体是功率相加（截面本身与 $N$ 无关），仪器透过率与探测效率也不在本节讨论范围内，
定量比较实验时须另行补齐。与量子理论一节「与入射光强度成正比、与距离平方成反比」两条
一致——后者说的是固定接收面积的情形，即上述第 3 条。

\emph{核对说明}：已对照原页面图像（书内 p.347 / PDF p.366）。原书在「振动的感生电偶极
引起的散射强度为」下印的是\textbf{完整式}
$I=k'_{\omega}\,\omega_{\mathrm{s}}^{4}p^{2}\sin^{2}\theta$，并给出前因子的显式定义
$k'_{\omega}=\dfrac{1}{32\pi^{2}\varepsilon_{0}c_{0}^{3}}$（$\varepsilon_{0}$、$c_{0}$
分别为真空介电常数与真空中的光速，书内 p.347 原印）。可见原书该式\textbf{本来就含频率
四次方}（$\omega_{\mathrm{s}}^{4}$，$\omega_{\mathrm{s}}$ 为散射角频率）与感生偶极矩平方
$p^{2}$，前因子是写成显式的常数而非未定系数；此前源 OCR 残片作
「$I=I_{0}\sigma\sin^{2}\theta$」（OCR 层失真），据此产生的「$\sigma$ 是否含 $\nu_{0}^{4}$
无法确认」一疑现予撤销。本稿把原书前因子 $k'_{\omega}\omega_{\mathrm{s}}^{4}p^{2}$
整体认作偶极垂直方向的参考方向系数 $K$ 并改写为辐照度口径的
式~\eqref{eq:classical-intensity}（$P_{\mathrm{sc}}=E_{0}\,(\mathrm{d}\sigma/\mathrm{d}\Omega)
\,\Delta\Omega$，$E_{0}=P_{0}/A_{\mathrm{beam}}$），角因子 $\sin^{2}\theta$ 只经
式~\eqref{eq:dipole-dcs} 的微分散面出现一次，不再在功率式中重复相乘；该改写只做口径统一，
$\nu_{0}^{4}$ 依赖仍以原书的 $\omega_{\mathrm{s}}^{4}$ 为依据保留，文中四条缩放推论可自行
代入核对。

经典理论能解释拉曼与瑞利散射的频率，但\textbf{不能解释斯托克斯线与反斯托克斯线强度为
什么差几个数量级}——因为它把分子当成经典振子，而两类过程的初始布居数不同，这是纯量子
效应。

\subsubsection{拉曼散射的量子理论}

量子图像下，频率 $\nu_0$ 的入射单色光视为能量为 $h\nu_0$ 的光子束。光子与分子之间可
发生弹性与非弹性两类碰撞：弹性碰撞中光子不改变频率、只改变方向，对应瑞利散射；非弹性
碰撞中光子与分子交换能量，散射光子得到能量的过程对应频率增加的反斯托克斯拉曼散射，散
射光子失去能量的过程对应频率减小的斯托克斯拉曼散射。

能级图像（图~\ref{fig:raman-levels}）是这一节的读图要点。处于基态 $E_0$ 的分子被
$h\nu_0$ 激发到\textbf{受激虚态}；虚态不稳定，电子很快跃迁回去。落回 $E_0$ 并放出
$h\nu_0$ 光子即瑞利散射；落到电子基态中的振动激发态 $E_n$ 并放出
$h(\nu_0-\nu_{\mathrm{m}})$ 光子即斯托克斯线；若分子原先已在 $E_n$，跃迁回 $E_0$ 时放出
$h(\nu_0+\nu_{\mathrm{m}})$ 光子即反斯托克斯线，此时分子失去 $h\nu_{\mathrm{m}}$ 的能量。

\begin{figure}[htbp]
\centering
\begin{tikzpicture}[
  x=1mm, y=1mm,
  font=\fontsize{8.5}{10}\selectfont,
  lvl/.style={line width=0.9pt},
  ray/.style={-{Latex[length=1.6mm]}, line width=0.7pt},
  stk/.style={-{Latex[length=1.6mm]}, line width=0.7pt, dashed}
]
% ---- 四个能级 ----
\draw[lvl] (5,0) -- (33,0)    node[right=1mm] {$E_0$（电子基态）};
\draw[lvl] (5,7) -- (33,7)    node[right=1mm] {$E_n$（振动激发态，$h\nu_{\mathrm{m}}$）};
\draw[lvl] (5,30) -- (33,30)  node[right=1mm] {$E_{\mathrm{v}}$：受激虚态（不稳定）};
\draw[lvl, dashed] (18,-4) -- (18,37) node[above=0.8mm, inner sep=0pt] {共振拉曼};
% ---- 吸收与发射 ----
\draw[ray] (7,0)   -- (7,30);
\draw[ray] (7,30)  -- (11,0)   node[midway, left=0.8mm, font=\fontsize{8}{9.6}\selectfont] {$\nu_0$};
\draw[ray] (13,0)  -- (13,30);
\draw[ray] (13,30) -- (17,7)   node[midway, left=0.8mm, font=\fontsize{8}{9.6}\selectfont] {$\nu_0-\nu_{\mathrm{m}}$};
\draw[stk] (21,7)  -- (21,30);
\draw[stk] (21,30) -- (25,0)   node[midway, left=0.8mm, font=\fontsize{8}{9.6}\selectfont] {$\nu_0+\nu_{\mathrm{m}}$};
\draw[ray] (27,2)  -- (27,30);
\draw[ray] (27,30) -- (31,2)   node[midway, left=0.8mm, font=\fontsize{8}{9.6}\selectfont] {$\nu_0'$};
% ---- 标注 ----
% 该 node 的文本会撑开整个 tikzpicture：必须给 text width，
% 否则最长的说明行直接决定图宽并顶出版心（本项目实测超 66 pt）。
\node[font=\fontsize{8}{9.6}\selectfont, align=left, anchor=north west,
        text width=118mm, inner sep=0pt] at (5,-11)
  {瑞利散射：$\nu_0$ 不变，强度最大，须被抑制\\[1pt]
   斯托克斯：起始布居在 $E_0$，强度大（\textbf{常规测量用它}）\\[1pt]
   反斯托克斯：起始布居在 $E_n$，强度小数个量级\\[1pt]
   共振拉曼：$\nu_0$ 接近真实电子跃迁能级，虚态不再是「虚」的，
   特定生色团信号被选择性增强};
\end{tikzpicture}
\caption{拉曼与瑞利散射的能级过程示意（含共振拉曼）。读者应看三件事：基态 $E_0$ 与振动
激发态 $E_n$ 之间相差一个振动量子 $h\nu_{\mathrm{m}}$，而\textbf{中间态是虚态、不是真实
能级}；斯托克斯线起始于布居数大的 $E_0$，反斯托克斯线起始于布居数小的 $E_n$，这是两者
强度相差数量级的几何来源；当激发频率接近真实电子跃迁时（右侧虚线处），共振拉曼使与生
色团相关的振动被选择性增强。局限：本图只画出一条振动模式，未画转动能级、热布居分布与
多声子过程。\schemamark\figsrc{据 S1（上册）图 6-1，书内 p.349 / PDF p.368 改绘。}}
\label{fig:raman-levels}
\end{figure}

\begin{principle}
\textbf{斯托克斯与反斯托克斯的强度比}。常规拉曼测的是斯托克斯线，其强度与反斯托克斯线
之比由两能级的布居数比决定。常温下按 Boltzmann 分布统计，处于基态的分子占绝大多数，
故斯托克斯线比反斯托克斯线强得多：
\begin{equation}
\begin{aligned}
\frac{I_{\mathrm{as}}}{I_{\mathrm{s}}}
&=\left(\frac{\nu_0+\nu_{\mathrm{m}}}{\nu_0-\nu_{\mathrm{m}}}\right)^{\!4}
\exp\!\left(-\frac{h c\,\Delta\tilde\nu}{k_{\mathrm{B}}T}\right),\\[1pt]
&\qquad \Delta\tilde\nu=\frac{\nu_{\mathrm{m}}}{c} .
\end{aligned}
\label{eq:stokes-ratio}
\end{equation}
\textbf{成立条件}：只有 Stokes 与 anti-Stokes 两条谱线来自同一振动模式、样品处于热平衡、
无共振增强、无自吸收；$\nu_0$ 与 $\nu_{\mathrm{m}}$ 是\textbf{频率}不是波数，指数里必须
用 $c\,\Delta\tilde\nu$ 把波数换算回频率。\textbf{量纲检查}：前因子为频率之比的四次方，
无量纲；指数中 $hc\Delta\tilde\nu$ 与 $k_{\mathrm{B}}T$ 同为能量量纲，故指数无量纲，两端
一致。\textbf{数值感觉}：$\Delta\tilde\nu=\qty{1000}{cm^{-1}}$、$T=\qty{298}{K}$ 时
$hc\Delta\tilde\nu/k_{\mathrm{B}}T\approx4.8$，指数因子约 \num{8e-3}，即反斯托克斯线
比斯托克斯线弱两个数量级以上。
\end{principle}

\emph{核对说明}：已对照原页面图像（书内 p.348 / PDF p.367，书内 p.350 / PDF p.369）。
式~\eqref{eq:stokes-ratio} 是按标准形式重建的，\textbf{原书未印出这条强度比公式}：原书在
该处只有定性表述，即「在常温下，根据玻尔兹曼分布统计，处于基态的分子占绝大多数，所以
通常斯托克斯线比反斯托克斯线强得多」（书内 p.348 原印）。原书量子理论一节另给反斯托克斯
强度式 $I^{\mathrm{AS}}_{\rho\sigma}=(k^{2}/R^{2})(\omega_{0}+\omega_{fi})^{4}
I^{i}_{\sigma_{0}}(\bar{\alpha}_{\rho\sigma})^{2}_{fi}$ 并明言「拉曼散射与入射频率的四次方
约成正比，与距离平方成反比，与极化率的导数的平方成正比，与入射强度成正比」（书内 p.350
原印），本式前因子的频率四次方即据此保留；指数取
$\exp(-h\nu_{\mathrm{m}}/k_{\mathrm{B}}T)$ 这一纯频率形式是标准写法，原书未印。
\textbf{数值算例}（$hc\Delta\tilde\nu/k_{\mathrm{B}}T\approx4.8$、约 \num{8e-3}）系本分片
依标准常数补算，属改编补充，不作为原书数据引用。

对拉曼散射强度，量子理论给出与经典结论一致的四条依赖关系：\textbf{与散射频率的四次方
约成正比、与距离平方成反比、与极化率导数的平方成正比、与入射光强度成正比}
\src{S1 §6.1.2，书内 p.350 / PDF p.369}。其中 $\nu_0^4$ 这一条在实验上最要紧：它意味着
用短波长激发可以显著提高散射强度，这正是紫外拉曼除了避开荧光之外还有灵敏度收益的原因。

\subsection{荧光的发生机制与抑制}

\subsubsection{荧光为什么能淹没拉曼信号}

分子的一个电子态由许多振动态组成，每个振动态又包含许多转动态，因此电子态包含大量精细
结构，构成一条准连续的能带。受光照射时，分子吸收光子能量，由电子基态跃迁到第一电子激
发态或更高激发态的振动和转动能级上；激发态分子不稳定，先通过无辐射跃迁降到第一电子激
发态的最低振转能级，再通过辐射跃迁回到基态并发射光子。由第一电子激发\textbf{单重态}所
产生的辐射跃迁伴随的发光称为\textbf{荧光}；由最低电子激发\textbf{三重态}所产生的辐射
跃迁伴随的发光称为\textbf{磷光}\src{S1 §6.1.3，书内 p.349—350 / PDF p.368—369}。

\begin{principle}
\textbf{两条信号的量级差是问题的全部要害}。拉曼是非弹性散射，每个分子只有约
\num{e-9}—\num{e-6} 的入射光子被拉曼散射（源文本给出的瑞利/拉曼强度比量级），不涉及
真实能级布居，没有增益；荧光是\textbf{吸收—弛豫—发射}过程，吸收截面大得多，且发射量子
产率可以接近 1。两者不在同一量级上：一旦激发线落进样品的电子吸收带，荧光可以比拉曼强
若干数量级，把微弱的拉曼峰整体抬到探测器的饱和区。更棘手的是，\textbf{荧光背景是一个宽
的连续谱}，它不像一条锐线那样可以与拉曼峰分辨——它直接抬高基线、压低相对对比度，使
信噪比按 $\propto\sqrt{t}$ 的缓慢规律改善，靠延长积分时间往往等不到收敛。
\end{principle}

\begin{dataflow}
从「谱图上有峰」到「这是拉曼峰」的判断链，在荧光存在时必须走完四步：
(1) \textbf{换激发波长重测}——真拉曼峰的\textbf{位移不动}，荧光带的位置随激发波长大幅
移动（因为荧光峰按波长标定、而拉曼按位移标定）；(2) \textbf{查强度比}——同一模式若在
多条激发线下都能测到而位移一致，才能排除荧光假峰；(3) \textbf{对基线}——宽荧光背景
必须扣除，扣基线的方式会直接改变弱峰的相对强度，凡做定量比较必须声明扣基线方法；
(4) \textbf{比空白}——测同条件下的纯载体与空样品池，扣除其中属于载体荧光与本底的那
一部分。缺任何一步，「\qty{980}{cm^{-1}} 处有一个峰，所以有表面钼物种」这类结论都可能
是荧光的产物。
\end{dataflow}

\subsubsection{抑制荧光的四条路线}

抑制荧光不是单一技巧，而是四条可并行的路线，各有代价（表~\ref{tab:fluor-suppress}）。

\begin{table}[htbp]
\centering
\caption{荧光抑制的四条路线及其代价。读者应看「代价与边界」一列：四条路线都不是免费的，
其中改变激发波长是首选，因为只有它在物理上避开了电子吸收带，其余三条都是与荧光做「减法」。
\src{据 S1 第 6 章引言，书内 p.345 / PDF p.364；§6.1.3—§6.2.1，书内 p.350—351 /
PDF p.369—370 归纳。}}
\label{tab:fluor-suppress}
{\srcsize
\begin{tabularx}{\columnwidth}{@{}lXX@{}}
\toprule
路线 & 做法与依据 & 代价与边界 \\
\midrule
改变激发波长 & 使激发线避开发色团的电子吸收带。紫外激发（如准分子激光）可把激发线移到吸收带短波侧；近红外激发则移到长波侧 & 紫外会同时增强共振效应，谱图不再「只反映浓度」；近红外散射强度按 $\nu_0^4$ 下降，信号更弱；换波长需要换光栅与探测器 \\
时间分辨 & 拉曼散射是瞬时的（约 \qty{e-14}{s}），荧光发射有纳秒级寿命，用锁模激光器加门控探测在荧光来得及发射之前或之外取信号 & 需要锁模激光器与门控（或条纹）探测系统，仪器复杂、成本高，且要求荧光寿命足够长 \\
傅里叶变换拉曼 & 用近红外（常见 \qty{1064}{nm}）激发加干涉仪，近红外激发本身大幅降低荧光量子产率，干涉仪的 Jaquinot 优势又提高光通量 & 近红外散射强度低、需大功率激光；水的吸收与样品自吸收、热效应需处理；低频截止受限 \\
猝灭 & 从源头减少发光物种：用高纯原料减少过渡金属离子（\ce{Fe^3+}、\ce{Cr^3+}、\ce{Mn^2+}）与有机物杂质；用氧化处理或强光照射分解有机杂质；通过焙烧减弱酸性位引起的荧光 & 属样品前处理，可能改变待测状态；猝灭不完全，对强荧光样品仍可能失效 \\
\bottomrule
\end{tabularx}}
\end{table}

\emph{核对说明}：已对照原页面图像（书内 p.345 / PDF p.364；书内 p.350—351 /
PDF p.369—370）。第 6
章引言写明「采用紫外激光作为激发光源的紫外拉曼光谱技术有效地解决了催化研究中所遇到的荧光干扰问题」，
即本表「改变激发波长」一路的核心依据；同时指出荧光干扰与灵敏度较低仍是拉曼光谱原位研究的难题。§6.1.3
荧光发生机制与 §6.2 仪器五部分组成见书内 p.350。§6.2.1
激光光源一节写明锁模激光器的成功应用使时间分辨光谱得以实现、激光全息技术的成功应用使傅里叶变换拉曼光谱得以实现，
分别对应本表「时间分辨」与「傅里叶变换拉曼」两路。本表「猝灭」一路（用高纯原料减少铁、铬、锰等过渡金属杂质，
焙烧减弱酸性位荧光）未在书内 p.345 与 p.350—351 出现，尚未逐页核对；
改变激发波长的两条具体波长走向（近红外移向长波侧、紫外移向短波侧）另见 §6.4.1（书内 p.374 / PDF
p.393），该页已核对。


\begin{pitfall}
三个最常见的误判。第一，\textbf{把扣基线后的谱峰当成原始证据}：宽荧光背景的扣除方式
不同，弱峰的相对强度可以相差数倍，凡是跨样品比较强度，必须统一并声明基线处理方法。
第二，\textbf{把荧光本身当作结构信息}：分子筛样品的荧光主要来自过渡金属离子、有机物杂
质和含铝分子筛的酸性位，它随原料与处理历史变化，与骨架结构没有一一对应关系。
第三，\textbf{认为「信噪比差就多测一会儿」}：荧光主导时延长积分时间主要累积背景，信噪
比改善极慢，此时正确做法是换路线（换激发波长、时间分辨或先做样品前处理），不是加时间。
\end{pitfall}

\begin{crosscheck}
\textbf{紫外拉曼是本章给出的第一条系统性答案}：采用紫外激光作为激发光源的紫外拉曼光谱
技术有效地解决了催化研究中所遇到的荧光干扰问题\src{S1 第 6 章引言，书内 p.345 /
PDF p.364}。它的机理与代价（共振增强带来的选择性、以及紫外光路与样品的特殊要求）在本
章后续节次展开。回查规则：凡在可见光激发下被荧光判定为「测不出来」的样品，应先问三
个问题——换更短的激发波长试过没有？样品是否已做去除有机杂质与过渡金属杂质的前处理？
能不能用时间分辨把纳秒级荧光与飞秒级拉曼分开？三者都试过仍失败，才可下「该样品无法用
拉曼表征」的结论。
\end{crosscheck}

\subsection{拉曼光谱实验技术的发展}

自 20 世纪 60 年代将激光器用于拉曼光谱仪后，仪器发展极快。一台激光拉曼光谱仪主要由
\textbf{激光光源、外光路系统、样品池、光谱仪、检测和记录系统}五部分组成
\src{S1 §6.2，书内 p.350 / PDF p.369}。本节按「部件—作用—对谱图质量的影响」压缩为
表~\ref{tab:raman-parts}，不逐类展开激光器与样品池的型号谱系。

\begin{table}[htbp]
\centering
\caption{激光拉曼光谱仪的五个部件：作用与对谱图质量的影响。读者应看最后一列——每个部件
都不是「可选配件」，它把一条具体条件写进了最终谱图；读别人的谱图时应能反推这些条件。
\src{据 S1 §6.2、§6.2.1—§6.2.5，书内 p.350—355 / PDF p.369—374 归纳。}}
\label{tab:raman-parts}
{\srcsize
\begin{tabularx}{\columnwidth}{@{}lXX@{}}
\toprule
部件 & 作用与主要选择 & 写进谱图的什么条件 \\
\midrule
激光光源 & 提供单色、高强度、方向性与偏振性好的激发线。可调谐激光器支撑共振拉曼；锁模激光器支撑时间分辨；常见激发线 \qty{1064}{nm}、\qty{532}{nm}、\qty{488}{nm}、\qty{355}{nm}、\qty{266}{nm} 等 & 决定拉曼位移的可用范围、散射强度（$\propto\nu_0^4$）、是否落入样品吸收带（荧光与共振）、退偏振测量是否可行 \\
外光路系统 & 激光器之后、单色器之前的光学系统：滤除气体激光的杂线与放电谱线、光斑整形、反射镜或棱镜改向、透镜聚焦到样品，再用聚光透镜把散射光聚焦到单色仪入射狭缝 & 决定杂散光水平、照射功率密度（是否烧样）、收集立体角（即信噪比）与偏振构型 \\
样品池 & 通光部分用熔融石英以便从紫外到近红外全波段激发，且杂质少、紫外下不发光；按实验要求做成不同形状。已有可旋转池（避免样品被打坏）、气氛可控可旋转的原位池、高温高压池与超高真空池；收集角常见 \qty{180}{\degree} 背散射、\qty{150}{\degree} 与 \qty{90}{\degree} 散射 & 决定可达到的温度/压力/气氛、样品是否被激光损伤、原位性与收集效率；收集角还影响偏振与定量比较 \\
光谱仪 & 色散型为主，要求杂散光小、色散度高。单联杂散光约 \num{e-4}，双联约 \num{e-8}，三联可达 \num{e-12}；常用 3—4 块光栅以覆盖紫外、可见、近红外 & 决定能测到多靠近瑞利线（低波数能力）、分辨率与弱峰可检性；多联光谱仪光通量下降，是灵敏度与杂散光的直接折中 \\
检测和记录系统 & 早期用光电倍增管配光子计数器或锁相放大器（要求高量子效率、低热离子暗电流）；近年在液氮冷却或风冷/水冷 CCD 上实现宽范围同时检测 & 决定积分时间、噪声来源（暗电流还是散粒噪声）、能否一次覆盖整段光谱；CCD 的宽范围同时检测大幅提高了信噪比并缩短总测量时间 \\
\bottomrule
\end{tabularx}}
\end{table}

\emph{核对说明}：已对照原页面图像（书内 p.350—351、p.355 / PDF p.369—370、p.374，
上册 §6.2、§6.2.1）。激光光源一行与原文相符：氩离子激光器最强谱线为 \qty{488}{nm} 与\qty{514.5}{nm}，
YAG 激光器波长为 \qty{1064}{nm}，准分子激光器中最具实用价值的是紫外准分子激光器，锁模激光器使时间分辨光谱得以实现（书内
p.351）。光谱仪一行与原文相符：单联杂散光约 \num{e-4}、双联约 \num{e-8}、三联可达 \num{e-12}，
常用 3 至 4 块光栅（书内 p.355）——本表原写「单联杂散光约 \num{e-5}—\num{e-4}」，其中下界
\num{e-5} 于原书无据，已按原书改为约 \num{e-4}。检测和记录系统一行与原文相符：早期用光电倍增管配光子计数器或锁相放大器，
近年在液氮冷却或风冷、水冷 CCD 上实现宽范围同时检测（书内 p.355）。尚待核对：外光路系统与样品池两行（通光部分用熔融石英、
\qty{180}{\degree} 背散射与 \qty{150}{\degree}、\qty{90}{\degree}
收集、可旋转池与原位池、高温高压池、超高真空池，见书内 p.352—354），以及表中 \qty{532}{nm}、
\qty{355}{nm}、\qty{266}{nm} 三条激发线，均尚未逐页核对。


\begin{measure}
选定一套条件前，按「先看荧光、再看杂散光、最后看灵敏度」的顺序做取舍，因为三者的代价
互相耦合：
\begin{itemize}
\item \textbf{激发波长}：先确认样品在可见区有没有强吸收。有荧光就优先向紫外或近红外
  移动；同时记住散射强度按 $\nu_0^4$ 变化，波长加倍信号就掉一个数量级。
\item \textbf{激光功率密度}：提高功率能线性提高拉曼强度，但热效应与光化学反应会改变样
  品状态，深色样品（如金属硫化物）尤其容易升温。可旋转样品池是常用对策。
\item \textbf{光谱仪联数}：需要低波数信息就必须用双联或三联压低瑞利线，代价是光通量
  下降；不需要低波数就不必付出这个代价。
\item \textbf{探测器}：宽范围同时检测优先选 CCD；需要极端弱信号的单点测量，光电倍增管
  配光子计数仍然可用。
\end{itemize}
本节的仪器条件是后续所有催化应用谱图的\textbf{解释前提}：不知道激发波长、功率密度和
光谱仪联数，就无法判断一张谱图里低波数区为什么是空的、弱峰是真峰还是噪声。
\end{measure}

\begin{methodsum}
\textbf{6.1—6.2 的方法骨架}。原理侧：拉曼测 $\partial\alpha/\partial q$，与红外的
$\partial\mu/\partial q$ 互补，有对称中心时互斥；拉曼位移
$\Delta\tilde\nu=(\nu_0-\nu_{\mathrm{s}})/c$ 是与激发波长无关的分子本征量；斯托克斯与反
斯托克斯的强度差来自 Boltzmann 布居。仪器侧：激光光源、外光路、样品池、光谱仪、检测
系统五个部件分别把激发条件、杂散光与收集效率、原位环境、低波数能力、噪声来源写进谱图。
\textbf{贯穿全章的单点问题仍然是荧光}——四条抑制路线中，改变激发波长（紫外拉曼）是
首选，时间分辨与傅里叶变换拉曼各有仪器代价，猝灭属样品前处理。
\end{methodsum}

\srcnote{本节据 S1（上册）第 6 章 §6.1—§6.2，书内 p.343—355 / PDF p.362—374 编写；
催化应用、紫外拉曼、表面增强拉曼与共焦成像见本章后续分片。文中未标注来源的公式推导步骤
（式 \eqref{eq:alpha-expansion} 的三角展开、式 \eqref{eq:stokes-ratio} 的数值算例）为本
分片作者补写，非原书内容。}
